% =========================
% report/notation.tex
% =========================

این بخش نمادها، تعاریف و مفاهیم ریاضی مورد استفاده را به‌صورت دقیق مشخص می‌کند.

\subsection{نمادها}
\begin{itemize}
  \item $\mathcal{I}_T$: تصویر الگو (Template)؛ می‌تواند نقشهٔ قطعه، شماتیک یا طرح مرجع باشد.
  \item $\mathcal{I}_Q$: تصویر پرسش/عکس واقعی (Query Photo) از قطعه در محیط واقعی.
  \item $\mathbf{x}=[x,y]^T$: یک نقطه دوبعدی در مختصات کارتزین (پیکسلی).
  \item $\tilde{\mathbf{x}}=[x,y,1]^T$: نمایش \textbf{همگن} نقطه؛ برای اعمال تبدیل‌های پرسپکتیو استفاده می‌شود.
  \item $\mathbf{H}\in\mathbb{R}^{3\times3}$: \textbf{هموگرافی}؛ نگاشت پرسپکتیو بین صفحهٔ الگو و صفحهٔ عکس.
  \item $\pi([u,v,w]^T)=[u/w,\, v/w]^T$: عملگر \textbf{پروجکشن} از مختصات همگن به کارتزین.
  \item $\tau$: آستانهٔ RANSAC برای خطای بازفرافکنی (واحد: پیکسل).
  \item $\alpha$: پارامتر هم‌پوشانی (Blend) با $0\le \alpha \le 1$.
\end{itemize}

\subsection{مفهوم مختصات همگن}
در هندسهٔ تصویر، تبدیل‌های پرسپکتیو را نمی‌توان با یک ماتریس $2\times2$ و بردار انتقال ساده توصیف کرد؛ به همین دلیل از مختصات همگن استفاده می‌شود. برای هر نقطه $\mathbf{x}=[x,y]^T$، نسخهٔ همگن $\tilde{\mathbf{x}}=[x,y,1]^T$ تعریف می‌شود. سپس هموگرافی اعمال می‌گردد:
\[
\tilde{\mathbf{x}}' = \mathbf{H}\tilde{\mathbf{x}} = [u,v,w]^T,
\quad \mathbf{x}' = \pi(\tilde{\mathbf{x}}')=[u/w,\, v/w]^T.
\]
بنابراین خروجی کارتزین با تقسیم بر مولفهٔ سوم ($w$) به دست می‌آید.

\subsection{اختصارات}
\begin{itemize}
  \item ORB: الگوریتم سریع استخراج ویژگی و توصیف‌گر دودویی \cite{Rublee2011ORB}.
  \item RANSAC: الگوریتم مقاوم برای برازش مدل در حضور داده‌های پرت \cite{FischlerBolles1981}.
  \item IoU: معیار اشتراک بر اجتماع برای ماسک‌ها (در صورت وجود زمین‌حقیقت).
\end{itemize}

\subsection{فرض کلیدی (اعتبار مدل)}
فرض اصلی این گزارش این است که ناحیهٔ مورد نظر روی قطعه یا الگو \textbf{تقریباً صفحه‌ای} است؛ در این حالت هموگرافی مدل مناسبی برای نگاشت بین دو تصویر است \cite{HartleyZisserman2004}. اگر سطح به‌شدت سه‌بعدی باشد یا خمیدگی قابل توجه داشته باشد، یک هموگرافی واحد ممکن است نتواند کل ناحیه را دقیق توصیف کند.
